\chapter[\hspace{0.6em}结构嵌套模型的 G 估计]{结构嵌套模型的 G 估计}

前两章我们用 IP 加权与标准化估计戒烟对体重增加的平均因果效应。本章介绍估计平均因果效应的第三种方法：g 估计。我们仍使用同一套 NHEFS 观测数据，并给出执行这些分析的简单计算机代码。

IP 加权、标准化与 g 估计常被合称为 g 方法，因为它们是为处理随时间变化的治疗所构成的广义治疗对比而设计的。本书第二部分中治疗不随时间变化，对这些情形套用 g 方法或许有些小题大做，因为另有更简单的途径。不过，在相对简单的设定中展示 g 方法，可以让我们专注于它们的主要特征，而不必应付第三部分所述的更复杂问题。

IP 加权与标准化在第一部分（第 2 章）引入，随后在第二部分配合模型加以描述（分别为第 12 章与第 13 章）。g 估计则一直留到第二部分才讲。原因在于：描述 g 估计时，先设定一个结构模型会方便得多，哪怕这个模型是饱和的。参数经由 g 估计得到的模型，称为结构嵌套模型。三种 g 方法建立在不同的建模假设之上。

\section{重访因果问题}

前两章我们用 IP 加权与标准化估计戒烟（治疗）$A$ 对体重增加（结局）$Y$ 的平均因果效应。为此，我们使用了 1566 名 25 至 74 岁吸烟者的数据：戒烟者归为受治疗 $A=1$，否则归为未受治疗 $A=0$。我们假定，在 $L$ 中各变量——性别、年龄、种族、教育、吸烟强度与吸烟年限、日常体力活动、休闲锻炼以及体重——的条件下，受治疗者与未受治疗者是可交换的。我们把差值尺度上的平均因果效应定义为 $\E[Y^{a=1,c=0}]-\E[Y^{a=0,c=0}]$，即若人人受治疗且无人被删失、与人人未受治疗且无人被删失这两种情形下本可观测到的平均体重之差。

\mnote{与前几章一样，我们把分析限制在性别、年龄、种族、体重、身高、教育、饮酒、基线（1971--1975）与随访（1982）时的吸烟强度与吸烟年限、以及基线时回答了病史问卷的 NHEFS 个体中。}

量 $\E[Y^{a=1,c=0}]-\E[Y^{a=0,c=0}]$ 度量的是整个总体中的平均因果效应。但有时人们关心的是总体某个子集中的平均因果效应。例如，可能想估计 45 岁个体的平均因果效应（$\E[Y^{a=1,c=0}\given age=45]-\E[Y^{a=0,c=0}\given age=45]$），或女性的、受教育程度较低者的，等等。要估计总体某一子集中的效应，可以使用含乘积项的边缘结构模型（见第 12 章），或只对该子集做标准化（第 13 章）。

设想研究者想估计戒烟 $A$ 对体重增加 $Y$ 的因果效应，而且是在由 $L$ 中各变量取值组合所定义的每一个层中分别估计。在我们的例子里，这样的层非常多。其中一层是｛未戒烟者、女性、白人、26 岁、大学辍学、每天 15 支烟、吸烟 12 年、中等锻炼、非常活跃、体重 112 kg｝。如第 4 章所述，拥有极大数据集的研究者可以把研究总体划分为互斥的子集或互不重叠的层，每一层由 $L$ 中变量的一组特定取值 $l$ 定义，然后估计每一层中的平均因果效应。12.5 节说明，另一种做法是把 $L$ 的全部变量、连同 $L$ 各分量与治疗 $A$ 之间的乘积项，一并加入边缘结构模型。此时稳定权重 $SW^A(L)$ 等于 1，不再需要 IP 加权，因为（未加权的）结局回归模型若设定正确，已充分校正了 $L$ 造成的全部混杂（见第 15 章）。

本章我们用 g 估计，在协变量 $L$ 定义的每一个层中估计戒烟 $A$ 对体重增加 $Y$ 的平均因果效应。这一条件效应表示为 $\E[Y^{a=1,c=0}\given L]-\E[Y^{a=0,c=0}\given L]$。在描述 g 估计之前，我们先介绍结构嵌套模型与秩保持，并在下一节用新的方式表述给定 $L$ 的可交换性条件。

\section{重访可交换性}

提醒一下（见第 2 章）：在我们的例子里，条件可交换性意味着，在研究总体中所有 $L$ 取值相同的个体所构成的任一子集内，那些未戒烟者（$A=0$）若戒烟，本会有与戒烟者（$A=1$）相同的平均体重增加；反之亦然。换句话说，条件可交换性意味着：若受治疗者与未受治疗者接受相同的治疗水平，两组的结局分布本会相同。

\mnote{本节第一段你也许会觉得重复而多余，因为我们此前已经讨论过条件可交换性的含义。若真如此，我们求之不得。}

当治疗水平 $a$ 下的结局 $Y^a$ 的分布在受治疗者与未受治疗者中相同时，各个反事实结局 $Y^a$ 都在协变量各水平内与实际的治疗水平 $A$ 独立，即对 $a=1$ 与 $a=0$ 都有 $Y^a\indep A\given L$。

取未治疗下的反事实结局 $Y^{a=0}$。当条件可交换性成立时，知道 $Y^{a=0}$ 的取值，无助于区分 $L$ 取某一特定值的戒烟者与未戒烟者。也就是说，对反事实结局 $Y^{a=0}$ 的一切取值，“是戒烟者”的（以 $L$ 为条件的）条件概率都相同。用数学语言写为

\[
\Pr[A=1\given Y^{a=0},L]=\Pr[A=1\given L]
\]

对二值治疗 $A$ 而言，这是条件可交换性的一种等价定义。

用治疗的条件概率来表述条件可交换性，对本章后面描述 g 估计会有帮助。

具体地说，设我们对治疗概率提出如下参数 logistic 模型：

\[
\operatorname{logit}\Pr[A=1\given Y^{a=0},L]=\alpha_0+\alpha_1Y^{a=0}+\alpha_2L
\]

\mnote{为简单起见，本书在认为不会引起混淆时，不区分向量参数与标量参数。}

其中 $\alpha_2$ 是参数向量，$L$ 的每个分量对应一个参数。若 $L$ 有 $p$ 个分量 $L_1,\dots,L_p$，则 $\alpha_2L=\sum_{j=1}^{p}\alpha_{2j}L_j$。这个模型与我们在第 12 章估计 IP 权重分母时所用的模型相同，只是它还把反事实结局 $Y^{a=0}$ 作为一个协变量纳入进来。

\section{结构嵌套均值模型}

我们想估计在 $L$ 各水平内治疗 $A$ 的平均因果效应，即
\[
\E[Y^{a=1}\given L]-\E[Y^{a=0}\given L].
\]
（为简单起见，先假定在本节稍后之前没有删失。）由于均值之差等于差值的均值，这一效应也可以表示为 $\E[Y^{a=1}-Y^{a=0}\given L]$。若 $L$ 不造成效应修饰，这些差值在各层间本应恒定，即 $\E[Y^{a=1}-Y^{a=0}\given L]=\beta_1$，此时 $\beta_1$ 既是每一层中的平均因果效应，也是整个总体中的平均因果效应。于是，关于条件因果效应的结构模型为 $\E[Y^a-Y^{a=0}\given L]=\beta_1a$。与关于条件均值 $\E[Y^a\given L]$ 的模型不同，关于均值差 $\E[Y^a-Y^{a=0}\given L]$ 的模型既不含截距 $\beta_0$，也不含 $\beta_2L$ 项，因为计算差值时这两项都抵消了。

更一般地，$L$ 可能造成效应修饰。例如，戒烟对重度吸烟者的因果效应也许大于轻度吸烟者。为允许因果效应依赖于 $L$，我们可以在结构模型中加入乘积项，即 $\E[Y^a-Y^{a=0}\given L]=\beta_1a+\beta_2aL$，其中 $\beta_2$ 是参数向量。在条件可交换性 $Y^a\indep A\given L$ 下，条件效应在受治疗者与未受治疗者中相同，因为在 $L$ 各水平内，受治疗者与未受治疗者平均而言是同一类人。因此，在可交换性下，结构模型也可以写成

\[
\E[Y^a-Y^{a=0}\given A=a,L]=\beta_1a+\beta_2aL
\]

这称为\emph{结构嵌套均值模型}。由 g 估计得到的参数 $\beta_1$ 与 $\beta_2$（后者同样是向量），量化了戒烟 $A$ 在 $A$ 与 $L$ 各水平内对 $Y$ 的平均因果效应。

第 13 章考虑了平均结局 $Y$ 的参数模型，那些模型同样以治疗 $A$ 与协变量 $L$ 为条件。在估计参数 g 公式时，那些结局模型是标准化的基础。与那些参数模型不同，结构嵌套模型是半参数的，因为它们对截距与 $L$ 的主效应都不作任何设定——也就是说，既没有参数 $\beta_0$，也没有 $\beta_3L$ 项的参数 $\beta_3$。由于把这些参数留作未指定，结构嵌套模型所作的假设更少，对模型错误设定也比参数 g 公式更稳健。关于结构嵌套模型与第 12 章边缘结构模型的关系，见精要 14.1。

\mnote{Robins (1994) 首次描述了结构嵌套模型这一族。当治疗随时间变化时，这些模型才是“嵌套”的。解释见第三部分。}

\begin{finepoint}[边缘结构模型与结构嵌套模型的关系]
考虑在连续协变量 $V$（$L$ 的一个分量）各水平内、治疗水平 $a$ 下平均结局的边缘结构均值模型
\[
\E[Y^a\given V]=\beta_0+\beta_1a+\beta_2aV+\beta_3V
\]
其和 $\beta_1+\beta_2v$ 是 $V=v$ 的个体中的平均因果效应 $\E[Y^{a=1}-Y^{a=0}\given V=v]$，而 $\beta_0+\beta_3v$ 则是这些个体在未治疗下的平均反事实结局 $\E[Y^{a=0}\given V=v]$。设唯一的推断目标是平均因果效应 $\beta_1+\beta_2v$，即我们无意估计 $\beta_0+\beta_3v=\E[Y^{a=0}\given V=v]$。那么我们就会把模型写成 $\E[Y^a\given V]=\E[Y^{a=0}\given V]+\beta_1a+\beta_2aV$，或等价地写成
\[
\E[Y^a-Y^{a=0}\given V]=\beta_1a+\beta_2aV
\]
这称为\emph{半参数边缘结构均值模型}，因为与第 12 章的边缘结构模型不同，它把未治疗下的平均反事实结局 $\E[Y^{a=0}\given V]$ 完全留作未指定。若只关心给定 $V$ 时 $A$ 的条件效应，那么当 $V$ 是连续变量或高维变量时，半参数边缘结构模型比参数模型更稳健：即使模型 $\beta_1a+\beta_2aV$ 正确，为 $\E[Y^{a=0}\given V]$ 设定的参数模型 $\beta_0+\beta_3V$ 若错误设定，也可能导致治疗效应的估计有偏。这一偏倚之所以产生，是因为 $(\beta_0,\beta_3)$ 的估计可能与 $(\beta_1,\beta_2)$ 的估计相关。

以可交换性所需的混杂因子 $L$ 的真子集 $V$ 为条件的半参数边缘结构模型，等同于给定协变量 $V$ 时处理瞬时效应（blip）的结构嵌套模型，例如 $\beta_1a+\beta_2aV$。因此，要在无删失时估计 $\beta_1$ 与 $\beta_2$，我们先以 IP 权重 $SW^A(V)=f(A\given V)/f(A\given L)$ 构造一个伪总体。在这个伪总体中只存在 $V$ 造成的混杂，因而该半参数边缘结构模型就是一个结构嵌套模型，其参数用 g 估计得到：把 $V$ 换成 $L$，并把每个个体的贡献以 $SW^A(V)$ 加权。

再考虑一个特例：半参数边缘结构均值模型以 $L$ 中\emph{全部}变量（而非仅子集 $V$）为条件，从而 $SW^A(V)$ 对所有个体都等于 1。也就是说，考虑模型 $\E[Y^a-Y^{a=0}\given L]=\beta_1a+\beta_2aL$，我们称之为\emph{仿半参数边缘结构模型}。在条件可交换性下，这个模型正是本章所用的结构嵌套均值模型。
\end{finepoint}

存在删失时，我们关心的因果效应不是 $\E[Y^{a=1}-Y^{a=0}\given A,L]$，而是若人人都保持未删失时的平均因果效应
\[
\E[Y^{a=1,c=0}-Y^{a=0,c=0}\given A,L].
\]
估计这一差值，需要同时校正治疗 $A$ 效应的混杂与（由删失 $C=1$ 造成的）选择偏倚。如前两章所述，IP 加权与标准化可用于校正这两种偏倚。而 g 估计只能用于校正混杂，不能校正选择偏倚。因此，使用 g 估计时，首先要用 IP 加权校正删失造成的选择偏倚。实践中，我们可以先估计删失的非稳定 IP 权重，构造一个无人被删失的伪总体，再对该伪总体应用 g 估计。

在我们的戒烟例子里，可以使用第 12 章估计的非稳定 IP 权重 $W^C=1/\Pr[C=0\given L,A]$。我们再次假定变量向量 $L$ 足以同时校正混杂与选择偏倚。

本章所述的一切 g 估计分析都纳入了 IP 权重，以校正删失可能带来的选择偏倚。在假定被删失者与未删失者在给定已测协变量 $L$ 时可交换的前提下，把结构嵌套均值模型 $\E[Y^a-Y^{a=0}\given A=a,L]=\beta_1a+\beta_2aL$ 应用于由 IP 权重 $W^C$ 构造的伪总体，实际上就是无删失情形下的结构模型：

\[
\E[Y^{a,c=0}-Y^{a=0,c=0}\given A=a,L]=\beta_1a+\beta_2aL
\]

\mnote{严格说来，当治疗是时固定的（而非随时间变化的）、不影响 $L$ 中任何变量、且结局只在单一时点测量时，用 g 估计校正选择偏倚并不需要 IP 加权。也就是说，若像我们一直假定的那样有 $Y^a\indep(A,C)\given L$，就可以对未删失个体直接应用 g 估计，而不必使用 IP 权重。}

为简单起见，本章以下省略上标 $c=0$。

本章我们用结构嵌套均值模型的 g 估计来估计二值治疗“戒烟”的效应，但结构嵌套模型也可用于连续治疗变量——如“吸烟强度的变化”（见第 12 章）。对连续变量，模型需要指定治疗 $A$ 对平均结局 $Y$ 之效应的剂量反应函数。例如 $\E[Y^a-Y^{a=0}\given A=a,L]=\beta_1a+\beta_2a^2+\beta_3aL+\beta_4a^2L$，或者 $\E[Y^a-Y^{a=0}\given A=a,L]$ 可以是 $A$ 与 $L$ 的光滑样条函数。关于二值结局的结构嵌套均值模型，见技术注记 14.1。

\begin{techpoint}[乘性结构嵌套均值模型]
正文只考虑加性结构嵌套均值模型。当结局变量 $Y$ 只能取正值时，人们往往更偏好乘性结构嵌套均值模型。乘性结构嵌套均值模型的一个例子是
\[
\log\frac{\E[Y^a\given A=a,L]}{\E[Y^{a=0}\given A=a,L]}=\beta_1a+\beta_2aL
\]
它可以按 14.5 节所述用 g 估计拟合，其中 $H(\psi^\dagger)$ 定义为 $Y\exp\{-\psi_1^\dagger a-\psi_2^\dagger aL\}$。

最初，上述乘性模型只能用于二值 $(0,1)$ 结局 $Y$，且要求 $Y=1$ 的概率在 $L$ 的所有层中都很小，以免模型预测出大于 1 的概率。Richardson、Robins 与 Wang (2017) 用条件对数优势积替代基线风险 $\Pr[Y=1\given A=0,L]$ 作为冗余参数，从而克服了这一稀有结局限制。这几位作者还把针对稀有二值结局的乘性结构嵌套均值模型推广到随时间变化的治疗，并用 g 估计构造因果参数的双重稳健估计量（Wang et al. 2022）。在这些进展之前，对非稀有的二值结局 $Y$，有人建议拟合结构嵌套 logistic 模型，如
\[
\operatorname{logit}\Pr[Y^a=1\given A=a,L]-\operatorname{logit}\Pr[Y^{a=0}=1\given A=a,L]=\beta_1a+\beta_2aL
\]
然而，结构嵌套 logistic 模型有两大缺点。第一，模型不可压缩，即边缘因果优势比不是各条件因果优势比的加权平均（精要 4.3）。第二，模型不易推广到随时间变化的治疗。细节见 Robins (1999) 与 Tchetgen Tchetgen and Rotnitzky (2011)。
\end{techpoint}

现在我们把注意力转向秩保持这一概念，它将帮助我们描述结构嵌套模型的 g 估计。

\section{秩保持}

在我们的戒烟例子里，所有个体都可以按观测结局 $Y$ 的取值排序。受试者 23522 体重增加 48.5 kg，排第一；个体 6928 体重增加 47.5 kg，排第二……个体 23321 体重增加 $-41.3$ kg，排最后。

\mnote{代码：程序 14.1}

同样，我们也可以设想按治疗下的反事实结局 $Y^{a=1}$ 的取值给所有个体排序——如果 $Y^{a=1}$ 的值对所有人都已知，而不只是对实际受治疗者已知。再设想我们真的能按 $Y^{a=1}$ 排序、也能按 $Y^{a=0}$ 排序。这样我们就有两份个体名单，各自按相应反事实结局从大到小排列。若两份名单的次序完全相同，我们就说存在\emph{秩保持}。

当治疗 $A$ 对结局 $Y$ 的效应在加性尺度上对研究总体中所有个体完全相同，我们说\emph{加性秩保持}成立。例如，若戒烟使每个人的体重都恰好增加 3 kg，那么按 $Y^{a=0}$ 排出的个体次序与按 $Y^{a=1}$ 排出的相同，只是后一份名单中每个人都要重 3 kg。加性秩保持的一个特例是严格因果零假设成立（见第 1 章），即治疗对研究总体中任何个体的结局都没有效应。就结构嵌套均值模型而言，我们关心的是 $L$ 各层内的加性秩保持。若治疗 $A$ 对结局 $Y$ 的效应对 $L$ 取值相同的所有个体完全相同，则\emph{条件加性秩保持}成立。

一个（加性条件）秩保持结构模型的例子是

\[
Y_i^a-Y_i^{a=0}=\psi_1a+\psi_2aL_i\qquad\text{对一切个体 }i
\]

其中 $\psi_1+\psi_2l$ 是对协变量取值 $L=l$ 的所有个体都相同的因果效应。也就是说，对每个 $L=l$ 的个体 $i$，$Y_i^{a=1}$ 等于 $Y_i^{a=0}+\psi_1+\psi_2l$。把个体在未治疗下的反事实结局 $Y_i^{a=0}$ 平移 $\psi_1+\psi_2l$，就得到她在治疗下反事实结局的取值。图 14.1 给出 $L=l$ 层内加性秩保持的一个例子。钟形曲线表示反事实结局 $Y^{a=0}$（左曲线）与 $Y^{a=1}$（右曲线）的分布。图上方两个点代表个体 $i$ 的两个反事实结局取值，下方两个点代表个体 $j$ 的两个反事实结局取值。箭头表示从 $Y^{a=0}$ 到 $Y^{a=1}$ 的平移，在该层内对所有个体都等于 $\psi_1+\psi_2l$。图 14.2 给出另一层 $L=l'$ 内秩保持的例子。该层反事实结局的分布与 $L=l$ 层不同。例如，图 14.1 中 $Y^{a=0}$ 的均值位于图 14.2 中 $Y^{a=0}$ 均值之左，这意味着平均而言，$L=l$ 层的个体在未戒烟时的体重增加小于 $L=l'$ 层的个体。对 $L=l'$ 的所有个体，从 $Y^{a=0}$ 到 $Y^{a=1}$ 的平移为 $\psi_1+\psi_2l'$，图中以个体 $p$ 与 $q$ 示出。

\mnote{图 14.1}

\begin{center}
\begin{tikzpicture}[
  x=1cm,y=2.7cm,
  declare function={g1(\x)=exp(-((\x-1.55)^2)/0.40); g2(\x)=exp(-((\x-2.95)^2)/0.40);},
  causal,font=\footnotesize]
  \draw[->] (0,0) -- (4.9,0) node[below right,inner sep=2pt] {体重增加};
  \draw[->] (0,0) -- (0,1.22) node[above,rotate=90,anchor=south,inner sep=1pt] {人数};
  \draw[semithick,domain=0.10:3.10,samples=90] plot (\x,{g1(\x)});
  \draw[semithick,domain=1.40:4.30,samples=90] plot (\x,{g2(\x)});
  % 个体 i：高位一对
  \draw[->] ({1.304},{g1(1.304)}) -- ({2.704},{g2(2.704)});
  \fill ({1.304},{g1(1.304)}) circle (1.9pt);
  \fill ({2.704},{g2(2.704)}) circle (1.9pt);
  \node[anchor=south east,yshift=2.5pt,inner sep=1pt] at ({1.304},{g1(1.304)}) {$Y_i^{a=0}$};
  \node[anchor=south east,yshift=2.5pt,inner sep=1pt] at ({2.704},{g2(2.704)}) {$Y_i^{a=1}$};
  % 个体 j：低位一对
  \draw[->] ({2.352},{g1(2.352)}) -- ({3.752},{g2(3.752)});
  \fill ({2.352},{g1(2.352)}) circle (1.9pt);
  \fill ({3.752},{g2(3.752)}) circle (1.9pt);
  \node[anchor=south east,yshift=2.5pt,inner sep=1pt] at ({2.352},{g1(2.352)}) {$Y_j^{a=0}$};
  \node[anchor=south east,yshift=2.5pt,inner sep=1pt] at ({3.752},{g2(3.752)}) {$Y_j^{a=1}$};
\end{tikzpicture}
\captionof{figure}{}
\label{fig:14.1}
\end{center}

\mnote{图 14.2}

\begin{center}
\begin{tikzpicture}[
  x=1cm,y=2.7cm,
  declare function={h1(\x)=exp(-((\x-1.70)^2)/0.30); h2(\x)=exp(-((\x-3.15)^2)/0.30);},
  causal,font=\footnotesize]
  \draw[->] (0,0) -- (4.9,0) node[below right,inner sep=2pt] {体重增加};
  \draw[->] (0,0) -- (0,1.22) node[above,rotate=90,anchor=south,inner sep=1pt] {人数};
  \draw[semithick,domain=0.35:3.05,samples=90] plot (\x,{h1(\x)});
  \draw[semithick,domain=1.80:4.50,samples=90] plot (\x,{h2(\x)});
  % 个体 p
  \draw[->] ({1.481},{h1(1.481)}) -- ({2.931},{h2(2.931)});
  \fill ({1.481},{h1(1.481)}) circle (1.9pt);
  \fill ({2.931},{h2(2.931)}) circle (1.9pt);
  \node[anchor=south east,yshift=2.5pt,inner sep=1pt] at ({1.481},{h1(1.481)}) {$Y_p^{a=0}$};
  \node[anchor=south east,yshift=2.5pt,inner sep=1pt] at ({2.931},{h2(2.931)}) {$Y_p^{a=1}$};
  % 个体 q
  \draw[->] ({2.261},{h1(2.261)}) -- ({3.711},{h2(3.711)});
  \fill ({2.261},{h1(2.261)}) circle (1.9pt);
  \fill ({3.711},{h2(3.711)}) circle (1.9pt);
  \node[anchor=south east,yshift=2.5pt,inner sep=1pt] at ({2.261},{h1(2.261)}) {$Y_q^{a=0}$};
  \node[anchor=south,yshift=3pt,inner sep=1pt] at ({3.711},{h2(3.711)}) {$Y_q^{a=1}$};
\end{tikzpicture}
\captionof{figure}{}
\label{fig:14.2}
\end{center}

对多数治疗与结局而言，不能指望个体因果效应在同协变量取值的个体之间恒定——甚至不能指望近似恒定——因此（加性条件）秩保持在科学上不可信。在我们的例子里，我们不认为戒烟对 $L$ 取值相同的所有人体重的影响都相同。即便在协变量 $L$ 各水平内，有些人在遗传上或其他方面也比另一些人更容易受戒烟影响。戒烟的个体因果效应因人而异：戒烟后有些人会增重很多，有些人增重很少，还有些人甚至会减轻体重。现实也许更接近图 14.3 描绘的情形：从 $Y^{a=0}$ 到 $Y^{a=1}$ 的平移在同协变量取值的个体间各不相同，连秩也不再保持，因为个体 $i$ 的结局在 $a=0$ 时小于个体 $j$ 的，在 $a=1$ 时却不然。

\mnote{图 14.3}

\begin{center}
\begin{tikzpicture}[
  x=1cm,y=2.7cm,
  declare function={g1(\x)=exp(-((\x-1.55)^2)/0.40); g2(\x)=exp(-((\x-2.95)^2)/0.40);},
  causal,font=\footnotesize]
  \draw[->] (0,0) -- (4.9,0) node[below right,inner sep=2pt] {体重增加};
  \draw[->] (0,0) -- (0,1.22) node[above,rotate=90,anchor=south,inner sep=1pt] {人数};
  \draw[semithick,domain=0.10:3.10,samples=90] plot (\x,{g1(\x)});
  \draw[semithick,domain=1.40:4.30,samples=90] plot (\x,{g2(\x)});
  \draw[->] ({1.345},{g1(1.345)}) -- ({3.155},{g2(3.155)});
  \fill ({1.345},{g1(1.345)}) circle (1.9pt);
  \fill ({3.155},{g2(3.155)}) circle (1.9pt);
  \node[anchor=south east,yshift=2.5pt,inner sep=1pt] at ({1.345},{g1(1.345)}) {$Y_i^{a=0}$};
  \node[anchor=south,yshift=3pt,inner sep=1pt] at ({3.155},{g2(3.155)}) {$Y_i^{a=1}$};
  \draw[->,dashed] ({2.352},{g1(2.352)}) -- ({3.548},{g2(3.548)});
  \fill ({2.352},{g1(2.352)}) circle (1.9pt);
  \fill ({3.548},{g2(3.548)}) circle (1.9pt);
  \node[anchor=north east,yshift=-2.5pt,inner sep=1pt] at ({2.352},{g1(2.352)}) {$Y_j^{a=0}$};
  \node[anchor=north west,yshift=-2.5pt,inner sep=1pt] at ({3.548},{g2(3.548)}) {$Y_j^{a=1}$};
\end{tikzpicture}
\captionof{figure}{}
\label{fig:14.3}
\end{center}

由于秩保持不可信，一般不应使用依赖它的因果推断方法。事实上，本书考虑的方法都不要求秩保持。例如，第 12 章的边缘结构均值模型是关于平均因果效应的模型，不是关于个体因果效应的模型，因而并不假定秩保持。戒烟对体重增加的平均因果效应估计为 3.5 kg（95\% 置信区间：2.5, 4.5）。这一平均效应与个体因果效应的秩保持是否成立无关。同样，上一节的结构嵌套均值模型也没有对秩保持作任何假定。

本节这个秩保持模型所作的假定，比非秩保持模型强得多：同一 $L$ 取值下所有个体的治疗效果恒定。没有理由在实践中使用这样一个不现实的秩保持模型。然而下一节我们仍用它来引入 g 估计，因为对秩保持模型来说 g 估计更容易理解，而且 g 估计过程对秩保持模型与非秩保持模型其实相同。注意，（条件加性）秩保持结构模型是一个结构均值模型——在 $L$ 各水平内，个体从 $Y^{a=0}$ 到 $Y^{a=1}$ 的平移的均值，等于每一个个体的平移。

\mnote{即使不存在秩保持，结构嵌套均值模型也定义良好。例如，对图 14.3 描绘的情形，即便同一 $L$ 取值下各个体的治疗效果并不完全相同，仍可以提出一个模型来估计 $L$ 各层内的平均因果效应。}

\section{G 估计}

本节把前三节的材料串联起来。设目标是估计结构嵌套均值模型 $\E[Y^a-Y^{a=0}\given A=a,L]=\beta_1a$ 的参数。为简单起见，先考虑只含一个参数 $\beta_1$ 的模型。由于模型没有乘积项 $\beta_2aL$，我们实际上假定了戒烟的平均因果效应在 $L$ 各层间不变，即 $L$ 不造成加性效应修饰。

我们还假定加性秩保持模型 $Y_i^a-Y_i^{a=0}=\psi_1a$ 对一切个体 $i$ 都正确设定。于是个体因果效应 $\psi_1$ 就等于我们所关心的平均因果效应 $\beta_1$。我们把秩保持模型写成 $Y^a-Y^{a=0}=\psi_1a$，不加下标 $i$ 来指个体，因为该模型对所有个体都相同。出于即将明了的原因，我们把这个模型写成等价形式

\[
Y^{a=0}=Y^a-\psi_1a
\]

g 估计的第一步，是把模型与观测数据联系起来。为此回想：由一致性，一个个体的观测结局 $Y$，在他受治疗 $A=1$ 时就是反事实结局 $Y^{a=1}$，在他未受治疗 $A=0$ 时就是反事实结局 $Y^{a=0}$。因此，若把结构模型中固定的取值 $a$ 换成每个个体各自的取值 $A$——有些人为 1，有些人为 0——就可以把反事实结局 $Y^a$ 换成该个体的观测结局 $Y^A=Y$。

秩保持结构模型于是给出一个方程，把每个个体的反事实结局 $Y^{a=0}$ 表示为其治疗与结局的观测数据以及未知参数 $\psi_1$ 的函数：

\[
Y^{a=0}=Y-\psi_1A
\]

若这个模型正确、且我们知道 $\psi_1$ 的值，就能算出研究总体中每个个体在未治疗下的反事实结局 $Y^{a=0}$。但我们并不知道 $\psi_1$。估计它正是分析的目标。

我们来玩个游戏。设想你的一位朋友知道 $\psi_1$ 的值，但他只告诉你 $\psi_1$ 是下列之一：$\psi^\dagger=-20$、$\psi^\dagger=0$ 或 $\psi^\dagger=10$。他向你挑战：“你能在这 3 个可能值 $\psi^\dagger$ 中认出真正的 $\psi_1$ 吗？”你接受挑战。对每个个体、每个候选值 $\psi^\dagger$，你计算

\[
H(\psi^\dagger)=Y-\psi^\dagger A
\]

新造出的变量 $H(-20)$、$H(0)$、$H(10)$ 是候选反事实。其中只有一个是反事实结局 $Y^{a=0}$。更确切地说，当 $\psi^\dagger=\psi_1$ 时 $H(\psi^\dagger)=Y^{a=0}$。在这个游戏里，选出正确的 $\psi_1$，等价于在三个候选反事实 $H(\psi^\dagger)$ 中挑出真正的反事实 $Y^{a=0}=H(\psi_1)$。你能想出挑选正确 $H(\psi^\dagger)$ 的办法吗？

回想 14.2 节：条件可交换性假设可以表述为给定反事实结局与协变量 $L$ 的治疗 logistic 模型。条件可交换性成立时，反事实结局的系数 $\alpha_1$ 应为零。于是我们有了一个简单办法，从三个变量 $H(\psi^\dagger)$ 中挑出真正的反事实：对每个候选 $H(\psi^\dagger)$ 各拟合一个 logistic 模型

\[
\operatorname{logit}\Pr[A=1\given H(\psi^\dagger),L]=\alpha_0+\alpha_1H(\psi^\dagger)+\alpha_2L,
\]

共三个。使 $\alpha_1=0$ 的那个候选 $H(\psi^\dagger)$ 就是反事实 $Y^{a=0}$，相应的 $\psi^\dagger$ 就是真值 $\psi_1$。例如，设 $H(\psi^\dagger=10)$ 在给定协变量 $L$ 时与治疗 $A$ 无关，那么 $\psi_1$ 的估计 $\hat\psi_1$ 就是 10。大功告成。这就是 g 估计。

\marginnote{\textcolor{black!85}{Rosenbaum (1987) 针对不随时间变化的治疗提出了这一过程的一个版本。}\par\vspace{4pt}\textcolor{black!85}{要点：g 估计并不检验条件可交换性是否成立；它假定条件可交换性成立。}}

然而在实践中，我们要在既没有知道答案、又不爱玩游戏的朋友的情况下 g 估计参数 $\psi_1$。因此需要在一切可能取值 $\psi^\dagger$ 上搜索，直到找到使 $H(\psi^\dagger)$ 的 $\alpha_1=0$ 的那个值。由于不可能检验所有取值——任何给定区间内都有无穷多个 $\psi^\dagger$——我们可以在 $-20$ 到 $20$ 的可能范围内以 $0.01$ 为步长做细搜索。搜索越细，得到的估计 $\hat\psi_1$ 越接近真值，计算量也越大。

\marginnote{\textcolor{black!85}{代码：程序 14.2}}[4.2\baselineskip]

在我们的戒烟例子里，我们先对 $\psi^\dagger$ 从 2.0 到 5.0、步长 0.1 的取值，计算每个个体在 31 个候选 $H(2.0)$、$H(2.1)$、$H(2.2)$、……$H(4.9)$、$H(5.0)$ 上的取值，然后拟合 31 个戒烟概率的 logistic 模型。这些模型与第 12 章估计 IP 权重分母所用的模型完全相同，只是每个模型都额外加入了一个候选 $H(\psi^\dagger)$。$H(\psi^\dagger)$ 的系数估计 $\hat\alpha_1$ 在 $H(3.4)$ 与 $H(3.5)$ 处最接近零。更细的搜索发现，$\hat\alpha_1$ 的最小值（实质上为零）出现在 $H(3.446)$。于是，戒烟对体重增加的平均因果效应 $\psi_1=\beta_1$ 的 g 估计为 $\hat\psi_1=3.4$ kg。

为计算 g 估计 3.4 附近的 95\% 置信区间，我们使用上述 logistic 模型中 $\alpha_1=0$ 的 Wald 检验 P 值。不出所料，当 $\psi^\dagger=3.446$ 时 P 值为 1——实际上是 0.998——因为这个 $\psi^\dagger$ 正是使候选 $H(\psi^\dagger)$ 的系数估计 $\hat\alpha_1=0$ 的取值。在我们对 $\psi^\dagger$ 从 2.0 到 5.0 拟合的 31 个 logistic 模型中，凡 $H(\psi^\dagger)$ 所依据的 $\psi^\dagger$ 大致在 2.5 到 4.5 之间的模型，P 值都大于 0.05。也就是说，用习惯的统计术语，对 2.5 到 4.5 之间这部分 $\psi^\dagger$ 取值，检验在 5\% 水平上“不拒绝零假设”。把检验结果反演，我们得出结论：3.4 附近 95\% 置信区间的上下限为 2.5 与 4.5。计算 95\% 置信区间的另一个办法，是对 g 估计过程做自助法。

\mnote{除 Wald 检验外，任何有效的检验都可以用。例如，得分检验能简化计算（无需拟合多个模型），且在大样本下与 Wald 检验实质上等价。}

更一般地，g 估计的 95\% 置信区间由这样一组 $\psi^\dagger$ 取值确定：检验 $\alpha_1=0$ 时 P 值大于 0.05 的那些值。95\% 置信区间通过对 $\alpha_1=0$ 的统计检验作反演得到，其上下限就是使 P 值大于 0.05 的 $\psi^\dagger$ 取值集合的边界。在我们的例子里，由于用 IP 加权校正删失，统计检验基于稳健方差估计量。因此我们的 95\% 置信区间在大样本下是保守的，即至少有 95\% 的概率覆盖真值。在大样本下，自助法会给出非保守、因而可能更窄的 g 估计 95\% 置信区间。

\mnote{存在删失时，logistic 模型的拟合必然限于未删失个体（$C=0$），且每个个体的贡献以其 IP 权重 $SW^C$ 的估计值加权。见技术注记 14.2。}

回到非秩保持模型。对 $\psi_1$ 的 g 估计算法（即实现该过程的计算机代码）给出均值模型参数 $\beta_1$ 的一致估计，前提是均值模型正确设定（即平均治疗效果在 $L$ 各水平上都相同）。这一点与个体治疗效果是否恒定无关，即与条件加性秩保持是否成立无关。换言之，g 估计算法的有效性并不真的要求对一切个体都有 $H(\beta_1)=Y^{a=0}$，其中 $\beta_1$ 是均值模型中的参数值。算法只要求 $H(\beta_1)$ 与 $Y^{a=0}$ 在给定 $L$ 时有相同的条件均值。

有趣的是，上述 g 估计过程可以很方便地加以修改，以纳入对未测混杂的敏感性分析，如精要 14.2 所述。

\begin{finepoint}[未测混杂的敏感性分析]
g 估计依赖于这样一个事实：若给定 $L$ 的条件可交换性成立，则 $\alpha_1=0$。现在考虑条件可交换性不成立的设定。例如，设配偶吸烟的个体戒烟的概率较低，而配偶的吸烟状况又与体重增加 $Y$ 的某些重要决定因素相关，这些因素未包含在 $L$ 中。也就是说，存在配偶吸烟状况造成的未测混杂。此时 $L$ 中变量不足以实现受治疗者与未治疗者的可交换性，治疗 $A$ 与反事实 $Y^{a=0}$ 在给定 $L$ 时相关。即 $\alpha_1\neq0$，我们无法按正文所述应用 g 估计。

但 g 估计并不要求 $\alpha_1=0$。设由于配偶吸烟状况造成的未测混杂，$\alpha_1$ 的预期值为 0.1 而非 0。那么我们可以按正文所述应用 g 估计，只是要检验 $\alpha_1=0.1$ 而非 $\alpha_1=0$。g 估计并不要求给定 $L$ 的条件可交换性成立，而要求非可交换性的程度——$\alpha_1$ 的取值——已知。g 估计的这一性质可用于对未测混杂做敏感性分析。

若我们认为 $L$ 可能不足以校正混杂，就可以在一系列不同的未测混杂情形下重复 g 估计分析，这些情形由 $\alpha_1$ 的一个取值范围内的一系列值代表，并绘出各情形下的效应估计。这样的图显示我们的效应估计对不同方向、不同大小的未测混杂有多敏感。这种方法的一个实际问题是：如何在 $\alpha_1$ 尺度上量化未测混杂（0.1 算不算很大的未测混杂？）。Robins、Rotnitzky 与 Scharfstein (1999) 给出了用 g 估计做未测混杂敏感性分析的技术细节。
\end{finepoint}

\section{含两个或更多参数的结构嵌套模型}

到此为止，我们考虑的都是只含一个参数 $\beta_1$ 的结构嵌套均值模型。缺少乘积项 $\beta_2aL$，意味着我们认为戒烟的平均因果效应不随 $L$ 的各层变化。若 $L$ 的某些分量 $V$ 确实造成效应修饰，而我们没有加入乘积项 $\beta_2aV$，结构嵌套模型就会被错误设定——我们的因果推断也就错了。这与饱和边缘结构模型 $\E[Y^a]=\beta_0+\beta_1a$ 形成对照：即使 $V$ 造成效应修饰，未加入 $\beta_2aV$ 与 $\beta_3V$ 项也不会使该模型错误设定。不以 $V$ 为条件的边缘结构模型估计总体中的平均因果效应，而以 $V$ 为条件的模型估计 $V$ 各层内的平均因果效应。结构嵌套模型按定义估计协变量 $L$ 各层内的平均因果效应，而非总体中的平均因果效应。存在效应修饰时，若在结构嵌套模型中省略乘积项，一般会因模型错误设定而导致偏倚。

所幸，上一节所述的 g 估计过程可以推广到含乘积项的模型。例如，设我们认为戒烟的平均因果效应依赖于基线吸烟强度 $V$。于是可以考虑结构嵌套均值模型 $\E[Y^a-Y^{a=0}\given A=a,L]=\beta_1a+\beta_2aV$。由于结构模型有两个参数 $\beta_1$ 与 $\beta_2$，我们也需要在 $\Pr[A=1\given H(\beta^\dagger),L]$ 的 IP 加权 logistic 模型中放入两个参数，其中 $\beta^\dagger=(\beta_1^\dagger,\beta_2^\dagger)$，$H(\beta^\dagger)=Y-\beta_1^\dagger A-\beta_2^\dagger AV$。例如，可以拟合 logistic 模型

\[
\operatorname{logit}\Pr[A=1\given H(\beta^\dagger),L]=\alpha_0+\alpha_1H(\beta^\dagger)+\alpha_2H(\beta^\dagger)V+\alpha_3L,
\]

找出使 $H(\beta^\dagger)$ 在给定协变量 $L$ 时与治疗 $A$ 独立的 $\beta_1^\dagger$ 与 $\beta_2^\dagger$ 取值组合。也就是说，要搜索使 $\alpha_1$ 与 $\alpha_2$ 同时为零的 $\beta_1^\dagger$、$\beta_2^\dagger$ 组合。由于模型有两个参数，搜索须在二维空间中进行。因此，系统的暴力搜索会比上一节所述麻烦得多。

\mnote{第 12 章讨论过，边缘结构模型有一个可取的性质——零保持：当无平均因果效应的零假设为真时，模型绝不会被错误设定。结构嵌套模型同样保持零假设。相比之下，参数 g 公式对时固定治疗保持零假设，但在治疗随时间变化的情形下会失去这一性质（见第三部分）。}

然而，尽管我们是用参数搜索来引入 g 估计的，g 估计一般并不需要对参数的可能取值做搜索。事实上，对这里讨论的这类线性均值模型，估计值可以直接用公式算出，即估计量有闭式解。对非线性结构嵌套均值模型，不存在闭式估计量，但可以使用基于导数的标准最优化技术，如 Newton--Raphson，因为 g 估计可以看作对模型参数求解一个估计方程（细节见技术注记 14.2）。对某些用于生存分析的结构嵌套模型，则必须做搜索，因为估计方程对模型参数不可微（见第 17 章）。

在我们的戒烟例子里，g 估计为 $\hat\beta_1=2.86$、$\hat\beta_2=0.03$。相应的 95\% 置信区间最便于用自助法计算。

\mnote{代码：程序 14.3}

在更一般的情形下，我们会考虑允许戒烟的平均因果效应随 $L$ 中变量的一切层变化的模型。对二值治疗，非饱和线性模型

\[
\E[Y^a-Y^{a=0}]=\beta_1a+a\sum_{j=1}^{p}\beta_{2j}L_j
\]

有 $p+1$ 个参数 $\beta_1,\beta_{21},\dots,\beta_{2p}$，其中 $\beta_{2j}$ 是乘积项 $aL_j$ 对应的参数，$L_j$ 代表 $L$ 的 $p$ 个分量之一。整个研究总体中的平均因果效应便可算作
\[
\beta_1+\frac{1}{n}\sum_{i=1}^{n}\sum_{j=1}^{p}\beta_{2j}L_{ij},
\]
其中 $n$ 是研究中的个体数。

介绍完 g 方法，下面我们回顾校正混杂的两种方法——结局回归与倾向性评分，它们可以说是最常用的。

\begin{techpoint}[结构嵌套均值模型的 g 估计]
考虑结构嵌套均值模型
\[
\E\bigl[Y-Y^{a=0}\bigm|A,L\bigr]=A\,\gamma(L;\beta)
\]
其中 $\gamma(L;\beta^\dagger)$ 是已知函数，$\beta^\dagger$ 通常是向量参数，且 $\gamma(L;\beta^\dagger)=0$ 当 $\beta^\dagger=0$。在正文所述各项假设（包括 $\E[A\given L]$ 的参数模型正确设定）下，用 g 估计可以得到 $\beta$ 的渐近无偏且渐近正态的估计。具体而言，我们的 $\beta$ 估计，是使 $H(\beta^\dagger)=Y-A\gamma(L;\beta^\dagger)$ 与 $A$ 在给定 $L$ 时的关联最小的那个 $\beta^\dagger$ 取值。当我们以得分检验（例如见 Casella and Berger 2002）为基础构造 g 估计时，这一过程等价于找出使下列估计方程成立的参数值 $\beta^\dagger$：
\[
\sum_{i=1}^{n}I[C_i=0]\,W_i^C\,H_i(\beta^\dagger)\bigl(A_i-\E[A\given L_i]\bigr)q(L_i)=0
\]
其中 $q(L_i)$ 是维数与 $\beta$ 相同的（由用户指定的）向量函数，$I[C_i=0]$ 是个体 $i$ 未删失的指示变量，IP 权重 $W_i^C$ 与期望 $\E[A\given L_i]=\Pr[A=1\given L_i]$ 都换成其估计值。$\E[A\given L_i]$ 可以用治疗对协变量 $L$ 的 logistic 模型估计，其中个体 $i$ 的贡献在 $C_i=0$ 时以 $W_i^C$ 加权，否则为零。［由于 $A$ 与 $L$ 在所有个体上都能观测到，我们也可以用全部个体对 $A$ 关于 $L$ 做未加权 logistic 回归来估计 $\E[A\given L_i]$。］向量函数 $q(L_i)$ 的选择影响估计量的统计效率，但不影响其一致性。也就是说，尽管 $q$ 的任何取法都能给出有效的置信区间，置信区间的长度却取决于该函数。Robins (1994) 对结构嵌套均值模型作了形式化描述，并导出了使置信区间长度最小的函数。

当 $\gamma(L;\beta^\dagger)$ 关于 $\beta^\dagger$ 线性，即 $\gamma(L;\beta^\dagger)=\beta^{\dagger,\mathrm T}d(L)$，其中 $d(L)$ 是维数与 $\beta$ 相同的已知向量函数时，该方程的解有闭式。此时若取 $q(L)=d(L)$，$\hat\beta$ 等于
\begin{multline*}
\Biggl(\sum_{i=1}^{n}I[C_i=0]\,W_i^C A_i\bigl(A_i-\E[A\given L_i]\bigr)d(L_i)d(L_i)^{\mathrm T}\Biggr)^{-1}\\
\times\sum_{i=1}^{n}I[C_i=0]\,W_i^C Y_i\bigl(A_i-\E[A\given L_i]\bigr)d(L_i)
\end{multline*}

一个自然而然的问题是：能否把上述估计方程中的 $H_i(\beta^\dagger)$ 换成非线性函数，例如 $[H_i(\beta^\dagger)]^3$，从而提高统计效率，同时仍保持估计的一致性？对于像本章所述结构嵌套均值模型这样、为识别只施加给定 $L$ 的均值独立条件（即 $\E[H(\beta_1)\given A,L]=\E[H(\beta_1)\given L]$）的模型，$H_i(\beta^\dagger)$ 的非线性函数不能用于我们的估计方程。$H_i(\beta^\dagger)$ 的非线性函数可用于施加分布独立条件（即 $H(\beta_1)\indep A\given L$）的模型，如结构嵌套分布模型（本章不介绍）：它把给定 $(A=a,L)$ 时 $Y^a$ 分布的各分位数，映射为给定 $(A=a,L)$ 时 $Y^0$ 分布的各分位数。

$\beta$ 的估计量只有在用于估计 $\E[A\given L]$ 与 $\Pr[C=1\given A,L]$ 的两个模型都正确时才一致。我们可以构造更稳健的估计量：把估计方程中的 $H(\beta^\dagger)$ 换成 $H(\beta^\dagger)-\E[H(\beta^\dagger)\given L]$，再用未删失个体拟合 $\E[H(\beta^\dagger)\given L]=\E[Y^{a=0}\given L]$ 的未加权线性模型，以估计后一个条件期望。若这个模型正确，那么即便上述 $\E[A\given L]$ 与 $\Pr[C=1\given A,L]$ 的模型都不正确，求解修改后估计方程所得的 $\beta$ 估计仍然一致。于是，只要 (i) $\E[H(\beta^\dagger)\given L]$ 的模型正确，或 (ii) $\E[A\given L]$ 与 $\Pr[C=1\given A,L]$ 的两个模型都正确——无需知道 (i) 与 (ii) 中哪一个正确——我们就得到 $\beta$ 的一致估计量。我们把这样的估计量称为\emph{双重稳健}的。技术注记 21.6 给出了治疗随时间变化的线性结构嵌套均值模型下该双重稳健估计量的闭式（见 Robins 2000）。
\end{techpoint}
